\documentclass[12pt]{article}
\usepackage{amsmath,graphicx}
\title{基于边缘计算的联邦学习优化方法研究}
\author{张三}
\date{2026年10月}

\begin{document}
\maketitle

% 这是注释，不应出现在检测结果中
\begin{abstract}
本文提出一种基于边缘计算的联邦学习优化方法。实验表明该方法在多个公开数据集上均取得了稳定的性能提升，同时将通信开销降低了约百分之三十。
\end{abstract}

\section{引言}\label{sec:intro}
近年来，随着人工智能技术的快速发展\cite{vaswani2017}，深度学习在自然语言处理领域取得了显著进展。然而在数据隐私敏感的场景下，集中式训练面临较大挑战。
背景相关内容。联邦学习允许各方在不共享原始数据的情况下协同训练模型\cite{mcmahan2017}。
相关工作表明，现有方法存在通信开销大、收敛速度慢等问题。文献~\cite{foo2020} 与 \cite{bar2021} 均讨论了这一问题，但均未考虑边缘节点的异构性。


\begin{figure}[htbp]
  \centering
  \includegraphics[width=0.8\textwidth]{arch.pdf}
  \caption{系统总体架构图。}
  \label{fig:arch}
\end{figure}

本文的主要贡献如下（系统总体架构如\autoref{fig:arch}所示）：
\begin{itemize}
  \item 第一项贡献是提出了新的梯度聚合策略。
  \item 第二项贡献是完整的实验评估与消融分析。
\end{itemize}

\begin{equation}
  \mathcal{L} = \sum_{i=1}^{n} \ell(f(x_i), y_i)
  \label{eq:loss}
\end{equation}

如式~\eqref{eq:loss}所示，损失函数由各样本损失累加而成，其中 $n$ 为样本数量，折扣因子 $\gamma\in(0,1)$ 控制收敛速度。

\section{结论}
如第~\ref{sec:intro}节所述，本文验证了所提方法的有效性，未来将探索更大规模的分布式场景。

\section*{参考文献}
\bibliographystyle{plain}
\bibliography{refs}

\appendix
\section{附录A：补充实验}
补充实验结果表明各组件均有效，超参数敏感性分析见下表。
\end{document}
